\chapter{Giới thiệu}
\label{chap:introduction}

\noindent
\textit{Chương này nêu lý do chọn đề tài, bài toán, mục tiêu, phạm vi và các đóng góp của đồ án, rồi tóm tắt cấu trúc
của báo cáo.}

\section{Lý do chọn đề tài}

Dung lượng của cell lithium-ion giảm dần trong quá trình sử dụng. Hệ thống quản lý pin (BMS) cần biết trạng thái sức
khoẻ (SOH) để ước tính quãng đường còn lại, điều chỉnh chiến lược sạc và cảnh báo thời điểm thay pin. Muốn biết chính
xác dung lượng phải xả đầy cell theo một quy trình chuẩn, mất nhiều giờ và phải dừng vận hành; ngược lại, điện áp và dòng
điện lúc sạc được ghi ở mọi chu kỳ của mọi cell. Vì vậy dữ liệu đầu vào thì nhiều, còn nhãn SOH thì ít.

Mạng nơ-ron có thông tin vật lý (PINN) là một cách tận dụng dữ liệu không nhãn: ràng buộc vật lý tính được chỉ từ dữ
liệu sạc. PINN4SOH~\cite{wang2024} áp dụng hướng này trên bốn bộ dữ liệu công khai và công bố mã nguồn. Tuy nhiên, khi
lọc ngoại lai và chuẩn hoá, bài dùng thống kê của cả vòng đời từng cell, kể cả cell test: đây là một dạng rò rỉ dữ liệu
đã biết~\cite{kapoor2023}, làm sai số đo được không phản ánh sai số khi triển khai; mức độ tái lập của số công bố cũng
chưa được kiểm tra. Đồ án
chọn đề tài này để xây dựng một mô hình ước lượng SOH khi ít nhãn và đánh giá nó, cùng các mô hình khác, trong một giao
thức không rò rỉ.

\section{Phát biểu bài toán}

Bài toán thứ nhất là ước lượng: cho 16 đặc trưng đoạn sạc và chỉ số chu kỳ của mọi chu kỳ, trong đó chỉ một phần các
cell huấn luyện có nhãn SOH, học một mô hình cho sai số nhỏ nhất trên các cell chưa từng xuất hiện khi huấn luyện; mô
hình chỉ được dùng dữ liệu tới chu kỳ hiện tại. Bài toán thứ hai là dự báo: đứng ở chu kỳ $k$, biết lịch sử SOH (đo
được hoặc ước lượng) tới $k$, dự báo SOH ở các chu kỳ sau đó, dung lượng còn lại và số chu kỳ còn lại tới ngưỡng cuối
đời (RUL). Phát biểu hình thức ở \Cref{sec:bai_toan}.

\section{Mục tiêu}

\begin{itemize}
  \item Xây dựng giao thức đánh giá không rò rỉ (chia theo cell, lặp lại, so ghép cặp có kiểm định) và đo tác động của
        các dạng rò rỉ thường gặp.
  \item Đề xuất mô hình PINN bán giám sát tận dụng cell không có nhãn, và mô hình kết hợp PINN-lai.
  \item So sánh có kiểm định với các mô hình đối chứng, gồm bốn mô hình học sâu của đề cương (MLP, CNN, LSTM,
        Transformer) và GBDT.
  \item So sánh với PINN4SOH trên đúng cách chia của bài, trong cùng điều kiện.
  \item Kiểm tra bằng số liệu vì sao dùng các đặc trưng đầu vào hiện có và vì sao bỏ các cách biến đổi khác.
  \item Mở rộng sang ngân sách đo dung lượng, cập nhật ước lượng khi dùng, dự báo SOH và RUL.
  \item Tổ chức mã nguồn sao cho mọi bảng, hình trong báo cáo sinh tự động từ kết quả và chạy lại cho kết quả trùng
        từng bit.
\end{itemize}

\section{Phạm vi}

Đồ án dùng bốn bộ dữ liệu công khai XJTU, TJU, MIT và HUST (387 cell, cell hình trụ 18650 sạc--xả trong phòng thí
nghiệm), ở bản đã trích 16 đặc trưng công bố kèm PINN4SOH. Đồ án không trích lại đặc trưng từ tín hiệu điện áp, dòng
điện gốc, không dùng dữ liệu vận hành ngoài phòng thí nghiệm hay dữ liệu nhiều cell ghép song song; các nội dung đó
thuộc giai đoạn 2 của đề cương. Mọi thí nghiệm chạy trên CPU.

\section{Đóng góp}

\begin{itemize}
  \item Giao thức đánh giá không rò rỉ cùng kiểm toán định lượng tác động của bốn dạng rò rỉ (\Cref{sec:protocol},
        \Cref{sec:leakres}).
  \item PINN bán giám sát với động học suy giảm dạng xám luôn không âm, và mô hình kết hợp PINN-lai
        (\Cref{sec:pinn,sec:pinnlai}), đánh giá ở bốn mức nhãn với chín mô hình đối chứng.
  \item So sánh với PINN4SOH trên đúng 39 nhóm đánh giá của bài, cùng cách chia và cách tính sai số
        (\Cref{ch:sosanh}).
  \item Phân tích đặc trưng đầu vào bằng tương quan hạng, bỏ nhóm, hoán vị và chọn trên validation
        (\Cref{ch:dactrung}).
  \item Đo dung lượng thưa trên mọi cell và cập nhật ước lượng khi cell mới được đo vài lần (\Cref{sec:mo_rong}).
  \item Dự báo SOH, dung lượng và RUL với hai mô hình có đặc trưng lịch sử: HiGBDT và HiPINN, trung vị của HiGBDT,
        PINN-dự báo và nhánh theo đời dựa trên đời còn lại dự đoán; thang đời, bộ chuẩn hoá và mọi lựa chọn chỉ dựa
        vào cell train và validation (\Cref{sec:e33pp}).
  \item Mã nguồn chạy được từ dữ liệu tới báo cáo bằng một lệnh, kiểm tra tái lập từng bit (\Cref{app:source-code}).
\end{itemize}

\section{Cấu trúc báo cáo}

Phần còn lại của báo cáo được tổ chức như sau. \Cref{chap:background} trình bày cơ sở lý thuyết và các nghiên cứu liên
quan. \Cref{chap:methodology} mô tả dữ liệu, đặc trưng đầu vào cùng số liệu chứng minh lựa chọn, giao thức đánh giá và
các mô hình. \Cref{chap:implementation} mô tả thiết kế và hiện thực hệ thống. \Cref{chap:evaluation} trình bày thiết
lập thực nghiệm, kết quả của đồ án và đối chiếu với các bài báo đã công bố, gồm cả các chỉ
số tốt hơn và các chỉ số chưa tốt. \Cref{chap:conclusion} kết luận. \Cref{app:source-code} cho mã nguồn và các kiểm
tra tái lập. Mọi con số trong bảng được sinh tự động từ tệp kết quả
của các lượt chạy.
